\section{Probability input, local geometry, and two arm consequences}
\label{sec:probability}

We separate the established percolation input from the consequences that
will be needed for rims.  All the input concerns the full-plane,
unit-intensity Poisson process with independent fair colors.  Write
$\Peta$ for the color law conditional on the point set $\eta$.

\subsection{The percolation input and its scope}

For $1\le r\le R$, let $A_j(r,R)$ be the arm event across the square
annulus $[-R,R]^2\setminus(-r,r)^2$.  For even $j$ the $j$ arms have
alternating colors in cyclic order.  For odd $j>1$ there are $j-1$
alternating arms and an additional black arm such that no Voronoi cell
intersects both that arm and any other arm.  The event $A_1$ is the existence
of a black arm.  In these definitions paths can be chosen so that their
used site sets are disjoint.  Set
\[
 \alpha_j(r,R)=\PP(A_j(r,R)),
 \qquad p_j^\eta(r,R)=\Peta(A_j(r,R)).
\]

\begin{proposition}[Established input]\label{prop:probability-input}
The following estimates hold at criticality.
\begin{enumerate}[label=(\roman*)]
\item For each fixed rectangle aspect ratio, both colors have crossing
probabilities bounded below by a positive constant, uniformly in its
size.  At sizes tending to infinity, the corresponding quenched crossing
probabilities are bounded below with probability tending to one.
More precisely, for each fixed finite family of buffered rectangle chains
on fixed square grids, and each $\gamma>0$, the required quenched crossing
lower bounds hold outside a geometry event of probability
$O_\gamma(R^{-\gamma})$ when the family is dilated by $R$.  The positive
lower bound may depend on the family and on $\gamma$.
\item For every fixed $j\ge1$, there is $C_j<\infty$ such that, for all
$1\le r\le R$,
\begin{equation}\label{prob:quenched-arms}
 \EE\bigl[p_j^\eta(r,R)^2\bigr]
       \le C_j\alpha_j(r,R)^2.
\end{equation}
\item There are constants $c_1,c_4>0$ and $C<\infty$ such that, for all
$1\le r\le R$,
\begin{equation}\label{prob:arm-powers}
 \alpha_1(r,R)\le C(r/R)^{c_1},\qquad
 \alpha_5(r,R)\le C(r/R)^2,\qquad
 \alpha_4(r,R)\le C(r/R)^{1+c_4}.
\end{equation}
\end{enumerate}
\end{proposition}

These estimates are due to
Tassion~\cite{Tassion2016}, Ahlberg--Griffiths--Morris--Tassion
\cite{AGMT2016}, and Vanneuville. The annealed scaling relations,
quasi-multiplicativity, and the five-arm estimate originate in
\cite{Vanneuville2019}. For the full-plane Poisson estimates above we use
\cite[Theorems~1.2--1.4]{Vanneuville2025}; the arbitrarily
polynomial exceptional-geometry estimate for the indicated grid families
is \cite[Proposition~3.14]{Vanneuville2025}, followed by finite gluing.
The powers in \eqref{prob:arm-powers} are, respectively,
\cite[(3.1), Proposition~3.4(iii), and Proposition~6.2]{Vanneuville2025}.
Only fixed relative grid meshes and a fixed finite list of buffered shapes
are used in applying Proposition~3.14.  Its minimum-size condition is then
met after dilation.  These assertions do not give crossing lower bounds
for arbitrary conditioned point processes.

We will use positive association and BK for increasing events conditional
on $\eta$.  We also need disjoint occurrence for events with prescribed
colors of both kinds.  For completeness, the required consequence of
Reimer's theorem has the following elementary reduction.

\begin{lemma}[Disjoint witnesses at fixed geometry]\label{prob:reimer}
For events $A,B$ in a finite product of fair bits, let $A\mathbin\square B$
mean that there are disjoint sets of coordinates certifying $A$ and $B$.
Then
\begin{equation}\label{prob:disjoint-product}
 \PP(A\mathbin\square B)\le\PP(A)\PP(B).
\end{equation}
In particular this inequality applies to bounded percolation tests
conditional on the Voronoi geometry, with disjoint used site sets.
\end{lemma}

\begin{proof}
Reimer's fair-bit inequality~\cite{Reimer2000} states that
$|A\mathbin\square B|\le |A\cap\overline B|$, where
$\overline B=\{\mathbf1-b:b\in B\}$ denotes bit complementation, not set
complementation.  Apply it on slices of pairs of configurations.
For each pair $(x,y)\in\{0,1\}^n\times\{0,1\}^n$, fix the set $K$ on
which they agree and their common values there.  On the other coordinates
the two strings are complementary.  If $x\in A\mathbin\square B$, its
disjoint certificates, restricted to the coordinates outside $K$, certify
the two restricted events on that slice.  Reimer's inequality bounds the
number of such pairs in the slice by the number with $x\in A$ and $y\in B$.
Summing over the slices gives
$2^n|A\mathbin\square B|\le |A||B|$, which is
\eqref{prob:disjoint-product}.  A fixed bounded region meets only finitely
many Voronoi cells almost surely, so the conditional percolation tests
indeed use a finite product.
\end{proof}

No annealed disjoint-occurrence product bound is asserted: averaging the
right side of \eqref{prob:disjoint-product} gives the expectation of a
product of quenched probabilities.  We will control that product using
\eqref{prob:quenched-arms}.

\subsection{Localization, moments, and buffered paths}

\begin{lemma}[Poisson mesh estimates]\label{lemma:mesh}
The following facts hold for the full-plane model of spacing $\eps$.
\begin{enumerate}[label=(\roman*)]
\item On a fixed compact set, the Delaunay triangulation and the Voronoi
cells meeting it are determined by the points in any fixed positive
collar, outside an event of probability tending to zero.  Their local
diameters tend to zero in probability.  The same determination statement
holds when the outside point configuration is changed while the collar
is retained.
\item For a unit square at spacing one, the probability that an incident
triangle has circumradius greater than $R>2$ is at most $Ce^{-cR^2}$.
The same form of bound holds for a Voronoi cell meeting the square whose
site has distance greater than $R$ from it.  Counts of incident triangles
and sums of any fixed polynomial diameter weights have moments of all
orders.
\item If $S$ is a square of side $s$, $k\ge0$, and $1\le p<\infty$, then
\begin{equation}\label{prob:mesh-moments}
 \left\|\sum_{t\cap S\ne\varnothing}\diam(t)^k\right\|_{L^p}
   \le C_{p,k}\eps^{k-2}(s+\eps)^2.
\end{equation}
For a deterministic Borel set $E$ of finite area,
\begin{equation}\label{prob:barycenter-intensity}
 \EE\sum_{z_t\in E}\diam(t)^k
    =m_k\eps^{k-2}\Area(E),\qquad m_k<\infty.
\end{equation}
At every fixed $s>0$, the sum of $\diam(t)^2$ over triangles with
barycenters in $S$ converges in every finite moment norm to $m_2s^2$ as
$\eps\downarrow0$.  The sum over triangles meeting $S$ has the same limit.
The corresponding statements hold for translation-covariant weights of
degree two bounded by a constant times $\diam(t)^2$, such as $|p_t|^2$
and $|q_t|^2$.
\item Within fixed positive margins, locally determined simple primal
cycles can approximate a fixed circle arbitrarily closely and in its
cyclic order.  Finitely many separated circles admit disjoint such cycles.
\end{enumerate}
\end{lemma}

\begin{proof}
If an empty disk of radius at least $\rho$ contains a point of a region
$K$, it contains an empty subdisk of radius comparable to $\rho$ with
center in an $O(\rho)$-collar of $K$.  A grid of mesh a small multiple of
$\rho$ reduces this to finitely many empty-disk tests.  The Poisson void
probability of each test is at most $\exp(-c\rho^2/\eps^2)$, and the
number of tests for a bounded region is at most a constant times
$(1+\diam(K)/\rho)^2$.

A triangle meeting $K$ has an empty circumdisk containing its meeting
point.  The same observation applies to a cell meeting $K$ with a remote
site.  Excluding the above empty subdisks makes all relevant circumdisks
small and places them inside the retained collar.  Both the positive
empty-circle tests defining those triangles and the exclusion of larger
ones are then local; they remain valid after any change outside that
collar.  This proves (i), including shrinking collar radii much larger
than $\eps\sqrt{\log(1/\eps)}$, and yields the exponential bounds in (ii).

On the event that all incident circumradii are at most $R$, the relevant
vertices lie in a box of side $O(R)$.  The incident triangle count is
bounded by a polynomial in the number of Poisson sites there, and every
diameter weight is bounded by a polynomial in $R$.  Splitting into
successive values of $R$, the exponential tail and the Poisson count
moments give all the asserted moments.  Tile $S$ by $O((s/\eps+1)^2)$
squares of side $\eps$.  Scaling the unit-square estimate and using the
triangle inequality in $L^p$ proves \eqref{prob:mesh-moments}; no
independence between adjacent squares is needed.  Stationarity and scaling
give \eqref{prob:barycenter-intensity}.

For a barycenter mask, the spatial ergodic theorem for the Poisson process,
applied at spacing one to expanding squares, proves the degree-two limit
after rescaling.  The higher moment bounds imply convergence in every
finite moment norm.  Triangles meeting the square but with barycenter
outside it are contained in a thin boundary neighborhood on a good mesh
event.  First fix the width of that neighborhood, use the same moment
bound there, and then let the width tend to zero; the bad mesh event is
removed by higher moments and its vanishing probability.  This proves the
statement for intersecting triangles and for the other indicated weights.

Finally, in a thin collar of a circle, follow successive triangles around
the circle and use their primal edges to form a walk with nonzero winding.
On a sufficiently fine mesh the walk stays in the collar and follows the
circle in order.  Decomposing this walk into simple cycles gives a simple
cycle with nonzero winding.  Choosing the collar within a prescribed
margin gives (iv), with simultaneous locality for any fixed finite list.
\end{proof}

At spacing $\eps$, write $A_j^\eps(r,R)$ for the corresponding physical
arm event, $p_{j,\eps}^\eta(r,R)$ for its quenched probability, and
$\alpha_j^\eps(r,R)=\alpha_j(r/\eps,R/\eps)$ for its annealed probability.
The arm estimates apply whenever $r\ge\eps$.  Thus every fixed positive
inner radius eventually satisfies their unit-scale lower cutoff.
Translation and rotation are allowed by
invariance of the Poisson law.  Tests in regions with disjoint positive
collars use independent point and color data after the local determination
of \cref{lemma:mesh}, up to a vanishing mesh exception.

Primal vertex paths and paths through the corresponding Voronoi cells
give the same buffered crossing tests.  For the direction from vertices
to cells, join sites to interior points of their shared cell sides, inside
the convex cells.  This gives an equivalent embedded drawing with the
same planar orders and vanishing displacement on a good mesh event.
Vertex-disjoint paths can be routed through just their own cells and the
open sides shared by successive cells; no other cells need be used or
intersected.  Conversely a cell crossing gives a connected chain of its
sites.  Fixed margins absorb the displacement in both directions.
Lower bounds in an open corridor are obtained by a finite chain of
overlapping fixed-shape rectangle crossings, with enough overlap and
margin to concatenate them.  Positive association conditional on $\eta$
and \cref{prop:probability-input} give a quenched positive lower bound
with high probability for each such fixed chain.

\subsection{A planar color-order fact}

Color-switching comparisons among prescribed arm colors are classical;
for triangular-lattice half-plane and annular versions see
\cite[Propositions~19--20]{Nolin2008}. The rectangular graph statement
needed here is proved directly below, so it remains valid after deleting
vertices and does not invoke a lattice theorem on Voronoi geometry.
We also use an elementary separation fact for triangulated disks.
Boundary sides below include their endpoints.  A path between opposite sides can always be
trimmed to be \emph{proper}, meeting those sides only at its first and
last vertices.

\begin{lemma}[Hex and indicator colorings]\label{prob:hex}
In a triangulated disk with four successive vertex sides, either there is
a black connection between the first and third sides or there is a white
connection between the second and fourth sides.  In particular, if a set
of vertices $X$ blocks every path between one pair of opposite sides,
then $X$ contains a path between the other pair.
\end{lemma}

\begin{proof}
Attach four exterior vertices, successively colored black, white, black,
white, and fan each onto its side; close the exterior to a sphere.  The
four color switches give four boundary ends of disagreement strands.
Their noncrossing pairing forces a connection between the two black fans
or between the two white fans.  Following the adjacent sites of a strand
gives the corresponding primal connection inside the disk.  For the last
assertion, color \emph{exactly} the vertices of $X$ black and every other
vertex white.  A white crossing of the blocked pair is impossible, so
the other crossing uses vertices of $X$ only.
\end{proof}

\begin{lemma}[Ordered colors in a rectangle]\label{prob:color-order}
Let a finite planar graph be drawn in a rectangle, with independent fair
colors on its vertices.  For any $k\ge1$, the probability of having $k$
vertex-disjoint crossings between its two end sides, with a prescribed
linear sequence of colors, depends only on $k$ and not on that sequence.
The assertion also holds for a planar subgraph with deleted vertices.
\end{lemma}

\begin{proof}
Fan the two end sides to two uncolored exterior terminals.  Crossings
become simple terminal-to-terminal paths, internally disjoint and ordered
from left to right.  Conversely, first-hit and last-hit trimming recovers
proper crossings in the same order.

Fix a first color.  If there is a path of that color, add an artificial
terminal-to-terminal edge on the exterior right.  The union of this edge
with all simple paths of the chosen color is biconnected: it is a union
of simple cycles sharing that edge.  Its left outer boundary is therefore
a simple path of the chosen color.  Call it the leftmost path.  It has
no same-colored bypass strictly on its left between distinct vertices,
including the terminals.  Conversely this property characterizes it.
Indeed a portion of any other such path lying strictly to its left,
between successive contacts, would be a bypass.  Uniqueness also follows
at the first divergence of two candidates, using the order at the start
terminal or the order from their common incoming edge.

The event that a specified path is leftmost is determined by its own
spins and the spins on its left.  All vertices strictly to its right
retain independent fair colors.  The leftmost path can replace the first
path of any ordered disjoint family: it lies on or to the left of that
path, while all the others lie strictly to its right.  Delete its internal
vertices and apply induction in the graph on its right.  For a fixed
first color this shows that the remaining color sequence does not affect
the probability.  A global color interchange changes the first color as
well.  The proof never uses the presence of every vertex of the original
rectangle, so also applies to a subgraph.
\end{proof}

\subsection{Four arms with both colors}

The next two estimates are local statements.  Let $x$ be deterministic and
$0<s<S$ fixed before $\eps$ tends to zero.  In a bulk patch assume a
neighborhood of $\overline{B(x,S)}$ is within the region of agreement with
the full-plane model.  All local comparisons below allow fixed changes
of the radii by positive constant factors.  The constants do not depend
on $x,s,S$, and bounded ratios $S/s$ are covered by increasing them.  The
statements make no assertion for a center chosen from the colors or for
an arbitrary microscopic choice $s=s_\eps$.

\begin{lemma}[Mixed-color four arms]\label{lemma:mixed-four}
For some $c>0$ and $C<\infty$,
\begin{equation}\label{prob:mixed-four}
 \limsup_{\eps\downarrow0}
 \PP\left(\begin{array}{c}
 \text{four vertex-disjoint primal arms across }B(x,S)\setminus B(x,s),\\
 \text{with both colors represented}
 \end{array}\right)
 \le C(s/S)^{1+c}.
\end{equation}
The four colors need not alternate.
\end{lemma}

\begin{proof}
We may take $S/s$ large.  Choose locally determined primal boundary
cycles strictly inside the annulus and trim the four witnesses to
disjoint proper paths between them.  Every step is made on a good mesh
event; its complement has vanishing probability at these fixed scales.
Select an adjacent pair of the four paths of opposite colors.  In their
gap there is a spanning disagreement strand $\gamma$: no strand can
leave through a monochromatic side path, and the odd number of switches
at each annular end forces a pairing from one end to the other.  Stop
$\gamma$ at its boundary edges.  The sites of its visited triangles are
disjoint from the two remaining paths, also at the ends because all four
paths are proper.

Sum over each fixed geometric candidate for $\gamma$ and each consistent
spin pattern on its visited sites.  Its occurrence is determined by
those spins.  Conditional on them, delete those sites and slit the annulus
along the strand.  The remaining sites are free fair bits in a planar
rectangle, and the two avoiding paths have a linear order.  By
\cref{prob:color-order}, the probability of the union of their four possible
color sequences is at most four times the probability of any specified
sequence.  Choose the sequence compatible with alternating four arms.
The two missing arms are supplied by the opposite-colored banks of
$\gamma$.  The sites beside each bank form a connected arm, are disjoint
from the two avoiding paths, and lie on their respective extreme sides
in the slit: they connect to the strand through its visited triangles
without crossing either avoiding path.  The resulting four arms thus
have the required cyclic order.

Let $N$ count proper spanning wall strands in the working annulus, each
stopped at its first boundary edges and counted once.  These are the
spanning components after the walls are cut at all contacts with the two
boundary cycles.  Let $G$ be the geometry event on which localization
holds and a fixed buffered subannulus has quenched white-circuit
probability at least $\delta>0$.  A fixed finite rectangle chain supplies
this circuit bound by \cref{prop:probability-input}.  In particular,
$\PP(G^c)=o(1)$ for each fixed $s,S$, and $\delta$ is independent of
$S/s$.  The preceding sum over strands gives
\begin{equation}\label{prob:strand-reduction}
 \PP(\text{mixed four arms})
 \le o(1)+C\EE\bigl[\ind_G N\ind_{A_4^\eps(r,R)}\bigr],
\end{equation}
where $r\asymp s$, $R\asymp S$ are deterministic buffered square radii.
To make this final transfer, first stop all the cell paths at their first
hit of the square outer boundary and then trim their starts at their last
hit of the inner boundary.  Planar order and disjoint used site sets are
preserved.  The summation is over all fixed strands; no selection of a
globally first strand is used when claiming that the remaining bits are
free.

On $G$, the variable $N$ has conditional moments of every fixed order,
uniformly in the annulus ratio.  Indeed its disjoint spanning strands
have adjacent black banks.  A gap between consecutive strands can
account for at most the two black banks bordering it; black paths in
different gaps cannot meet or cross a strand.  This also separates their
endpoint sites: a stopped strand ends inside a boundary edge between
opposite-colored sites, and each black boundary site belongs to only one
gap.  Thus $N\ge2k$ yields at least $k$ black crossings with disjoint
site witnesses, including their endpoints.  A black crossing of the
working annulus is incompatible with the white circuit in the fixed
subannulus, so its conditional probability
is at most $1-\delta$.  BK gives
\[
 \Peta(N\ge2k)\le(1-\delta)^k\quad\text{on }G,
 \qquad
 \EE[\ind_GN^q]\le C_{q,\delta}\quad(q<\infty).
\]
Keeping the indicator of $G$ in \eqref{prob:strand-reduction}, H\"older's
inequality and \eqref{prob:arm-powers} bound its second term by
\[
 C_q\PP(A_4^\eps(r,R))^{1-1/q}
 \le C_q(s/S)^{(1+c_4)(1-1/q)}.
\]
Choose $q$ large enough that the displayed exponent is greater than one,
and reduce $c$ accordingly.  Taking the microscopic limit proves the
claim.  No estimate of $N$ on bad geometry is required.
\end{proof}

\subsection{Two separated returns of a rim}

\begin{lemma}[Near-return bound]\label{lemma:near-return}
For some $c>0$ and $C<\infty$, and the same fixed-scale convention,
\begin{equation}\label{prob:near-return}
 \limsup_{\eps\downarrow0}
 \PP\left(\begin{array}{c}
 \text{some rim has two visits to }B(x,s)\text{ such that both}\\
 \text{complementary rim arcs between the visits leave }B(x,S)
 \end{array}\right)
 \le C(s/S)^{2+c}.
\end{equation}
The rim can be part of a global spherical completion.  The event is
bounded by a local monochromatic-arm test, without a union over rims.
\end{lemma}

\begin{proof}
Suppose first that the rim $P$ is black.  For a large enough ratio $S/s$,
take four nested primal cycles near radii
$2s,5s,S/5,S/2$.  They are locally determined and pairwise disjoint on
a good mesh event.  The two visits and the excursions in both rim
directions give four vertex-disjoint black journeys between the cycle near
$2s$ and the cycle near $S/2$.  Trim them properly and denote them by
$P_1,\ldots,P_4$ in their cyclic order.  The margins at both ends ensure
that the vertices between the middle cycles are interior vertices of
the original rim journeys, with both incident rim edges still present.
Retain these four complete journeys throughout the argument.

Each oriented journey has its right side in one of the two adjacent
annular gaps.  Since a gap has only two sides, at least two distinct gaps
have a right-facing rim side.  Fix one such gap.  The vertices of the
inner middle cycle block every connection between its annular ends.
Apply the indicator-coloring assertion of \cref{prob:hex}, coloring
exactly those vertices black.  It gives a transverse crossing between
the two journey sides, using only vertices of that cycle.  Trim this to
a proper crosscut.  Next work in the subdisk between this crosscut and
the outer end.  The outer middle cycle still separates the two ends,
because the first crosscut lies strictly inside it.  The same indicator
argument gives a second proper transverse crosscut using its vertices.
This successive construction orders the crosscuts even when the
original journeys make extra visits to either middle cycle.

Consider the quadrilateral between these crosscuts.  There is no actual
black crossing between its two journey sides.  Otherwise take a simple
such crossing and start just after its last contact with the right-facing
side.  Both incident rim edges at that departure vertex belong to the
original journey; the new edge therefore enters the open right sector of
the whole rim.  Follow the crossing only until its first subsequent
contact with the \emph{whole} rim, including any other rim portions in
the quadrilateral.  Such a contact exists because the other selected
side is part of the rim.  Simplicity makes the two contact vertices
distinct.  Until this contact the path stays in the right disk and has
no rim vertices in its interior, so it is a bypass forbidden by
\cref{lemma:no-bypass}.  This argument does not require the quadrilateral
itself to be contained in the smaller annulus between the middle cycles.
Hex in the actual coloring now gives a white crossing between the two
transverse cuts.  Trim it properly between the middle cycles.

Doing this in two distinct original gaps gives two disjoint white arms:
the boundaries shared by original gaps are black journey vertices, which
the white paths cannot use.  These white arms divide the middle annulus
into two interval quadrilaterals.  Arbitrarily trim the original black
journeys to obtain four disjoint black arms in this annulus.  It remains
to ensure that both white intervals contain a black arm; arbitrary trims
alone need not ensure this if a journey revisits a middle boundary.

For this purpose let $X$ be the set of vertices of all four \emph{original}
journeys, restricted to one of the white intervals.  Every path between
its two white sides must meet $X$.  Otherwise it would join two distinct
original gaps while avoiding all four paths that separate those gaps in
the larger working annulus.  Apply \cref{prob:hex} with the artificial
black set exactly $X$.  It gives a crossing between the middle cycles
using sites of $X$ and hence genuinely black sites.  The same argument
works in the other white interval.  If the four arbitrary black trims
already occupy both intervals, retain them.  If all lie in one, retain
three and add the newly obtained black crossing in the other.  The
white arms separate the intervals, so the four chosen black paths are
still vertex-disjoint.  The new crossing need not be a subpath of a
single original journey; membership of its sites in $X$ is enough.

We now have four black arms and two white arms, with at least one black
arm in each interval between the whites.  Choose one black arm in each
interval to obtain alternating four, add one of the two remaining black
arms as the extra arm for $A_5$, and reserve the last black arm for
$A_1$.  These witnesses use disjoint site sets.  Transfer all six paths
to deterministic buffered square boundaries, stopping at first outer
hits and last inner hits to preserve their order.  The event implies
$A_5^\eps(r,R)\mathbin\square A_1^\eps(r,R)$ for
$r\asymp s$, $R\asymp S$.

Conditional on $\eta$, \cref{prob:reimer} bounds its probability by
$p_{5,\eps}^\eta(r,R)p_{1,\eps}^\eta(r,R)$.  Averaging and applying
Cauchy--Schwarz
and \eqref{prob:quenched-arms} gives
\[
 \EE[p_{5,\eps}^\eta(r,R)p_{1,\eps}^\eta(r,R)]
 \le C\alpha_5^\eps(r,R)\alpha_1^\eps(r,R)
 \le C(s/S)^{2+c_1}.
\]
For a white rim interchange the colors, at a cost of a factor two.
Remove the mesh exception and take the limit superior.  The global rim
was used only for no-bypass; the resulting arm event is local, so the
same bound covers existence of any such rim in a bulk completion.
\end{proof}

Taking the smaller of the positive exponents in
\cref{lemma:mixed-four,lemma:near-return}, we henceforth use one constant
$c>0$ for both estimates, reducing it further when a later H\"older
inequality requires this.

\subsection{Facing arms and two-site gates}

A rim that visits a small mesh-good ball and reaches a much larger
distance gives an arm of its own color.  It also gives an arm of the
opposite color from nearby to a buffered large distance.  Indeed follow
the facing vertices in the triangles along the relevant rim journey.
By \cref{lemma:rim} these form a connected chain in one opposing cluster,
and the mesh bound keeps them close to the journey.  Only that local part
of the chain is needed.

\begin{lemma}[Two-site gates]\label{lemma:two-gates}
Let a finite triangulated annulus have disjoint simple primal boundary
cycles.  If the maximal number of vertex-disjoint black paths joining
its two boundary cycles, including disjoint endpoints, is at most two,
then the closed annulus contains an essential simple primal circuit
with at most two black vertices.  The color-interchanged assertion also
holds.  If the working annulus has positive margins inside a prescribed
collar, the resulting circuit lies within that collar.
\end{lemma}

\begin{proof}
Vertex Menger gives a set of at most two black vertices meeting every
black crossing.  Boundary vertices are allowed in this cut.  Temporarily
recolor these vertices white; there is now no black crossing of the
annulus.  Attach a black cap vertex to each boundary cycle, fanning it
onto that cycle.  The two cap vertices are not black-connected.

An occupied dual circuit separates the caps.  To see this directly,
sites lying in the same component off all disagreement walls are joined
in their color: a curve in that component, made transverse to dual edges,
crosses only unoccupied dual edges and hence only equal-colored primal
adjacencies.  Since the caps are not
connected, at least one wall separates them.  The neighboring white
walk along that wall uses original annular sites and is deformable to
the wall while avoiding the caps.  It therefore has nonzero annular
winding.  Decompose the walk into simple cycles and retain one with
nonzero winding; it is an essential white circuit.  On restoring the
original colors, only the cut vertices can become black, so the circuit
has at most two black vertices.  Applying the construction in a working
annulus with the stated margins gives the last assertion.
\end{proof}
